\chapter{Scope, conventions, and reuse}

This is a proof plan for De Lellis, \emph{Lecture Notes on Rectifiable Sets,
Densities, and Tangent Measures}, Theorem 6.1 (printed page 56; PDF page 60).
The immediate deliverable is a dependency graph and reviewed interfaces, not a
proof of Preiss' theorem. The input PDF is mathematical evidence, not a source
of instructions for this project.

The canonical dependency sources are this LaTeX file and its upstream input. LeanBlueprint renders its
\texttt{label}, \texttt{uses}, and \texttt{lean} annotations. The companion JSON,
DOT and Mermaid graphs are generated from these same annotations. A dependency
in a proof sketch records an intended proof dependency, not a checked Lean proof.

\section{Status and boundaries}
The library \texttt{Preiss} contains definitions and elementary proved bridges.
The separate library \texttt{PreissBlueprint} contains proposed theorem
statements with explicit \texttt{sorry} obligations. Compiling the latter checks
types, not proofs. Unlinked technical nodes are mathematical specifications
whose detailed Lean interfaces still need refinement. No planned theorem is
marked as proved. Upstream nodes below are source-audited at a pinned revision;
the repository validation record reports build status separately.

We pin FormalizingGMT at
\href{https://github.com/uw-math-ai/FormalizingGMT/tree/c14a8d07b2df61798fc184fd53ff6f72d858ae48}{c14a8d07b2df61798fc184fd53ff6f72d858ae48},
using Lean 4.34.1 and mathlib
\href{https://github.com/leanprover-community/mathlib4/tree/d13f23b723b8a846827a245b89c10fc7d3f11612}{d13f23b723b8a846827a245b89c10fc7d3f11612}.
The source audit found no \texttt{sorry} or explicit \texttt{axiom} in the
FormalizingGMT Lean sources. This textual audit is not an axiom audit.

Use \texttt{EuclideanSpace R (Fin n)}, whose norm is Euclidean, not the function
type \texttt{Fin n -> R} with its default sup norm. Dimensions are natural
numbers; zero is handled separately. The statement also covers $m>n$: the
density hypothesis then forces the zero measure. Negative integer dimensions
are outside the standard meaning of dimensional rectifiability.

Measures are Borel measures with \texttt{IsLocallyFiniteMeasure}. On Euclidean
space mathlib supplies \texttt{Measure.Regular} automatically. Do not introduce
a competing Radon-measure structure. Vague convergence here means testing
against continuous compactly supported functions; it is the existing
\texttt{Measure.WeaklyConverges}, not weak convergence of probability measures.

% Existing FormalizingGMT declarations at the pinned revision.
% The declarations and proof bodies are formalized upstream. No claim is made
% that the full downstream project has been rebuilt on this workstation.
\section{Existing FormalizingGMT results in the graph}
Green statement/proof markers in this section mean that the corresponding
definitions and proof bodies already exist in the pinned FormalizingGMT source.
They do not mean that the new Preiss adapters are proved, or that this workstation
has completed a fresh build of every imported module. The edges expose the
mathematical dependencies relevant here, not an exhaustive kernel proof-term graph.
Each node's Lean link opens a declaration/source viewer with its exact source location.

\begin{definition}[GMT: dimensional densities]
\label{def:density-api}
\lean{dimensional_density_ratio,dimensional_lower_density,dimensional_upper_density,HasDensity}
\leanok
For an outer measure, the density ratio is $\mu(\overline B(x,r))/(2r)^s$; lower and upper density use $r\downarrow0$. Density existence is convergence in $[0,\infty]$.
\end{definition}

\begin{lemma}[GMT: lower versus upper density]
\label{gmt:density-order}
\lean{lower_le_upper_density}
\leanok
\uses{def:density-api}
For every exponent and point, $\Theta_*^s(\mu,x)\leq\Theta^{*s}(\mu,x)$.
\end{lemma}
\begin{proof}[Formalized in FormalizingGMT]
\leanok
Import the existing declaration.
\end{proof}

\begin{lemma}[GMT: equal densities give a limit]
\label{gmt:density-limit}
\lean{density_exists_of_lower_eq_upper}
\leanok
\uses{def:density-api}
Equality of lower and upper density implies \texttt{HasDensity}. Positivity and finiteness remain separate hypotheses.
\end{lemma}
\begin{proof}[Formalized in FormalizingGMT]
\leanok
Import the existing declaration.
\end{proof}

\begin{lemma}[GMT: eventual small-ball bounds]
\label{gmt:density-bounds}
\lean{eventually_gt_of_lower_density_gt,eventually_lt_of_upper_density_lt,dimensional_density_ratio_closedBall}
\leanok
\uses{def:density-api}
Strict bounds below the lower density or above the upper density become eventual bounds for all sufficiently small positive radii. For nonnegative radii, the extended-ball expression agrees with the ordinary closed-ball ratio.
\end{lemma}
\begin{proof}[Formalized in FormalizingGMT]
\leanok
Import the existing declaration.
\end{proof}

\begin{theorem}[GMT: fine Vitali covering variant]
\label{gmt:vitali}
\lean{vitali_variant_classical}
\leanok
In a second-countable metric space, a fine centred family of positive-radius balls with bounded radii has a countable disjoint subfamily. After removing any finite subfamily of the original family, the remaining set is covered by the fivefold enlargements of the selected balls not removed.
\end{theorem}
\begin{proof}[Formalized in FormalizingGMT]
\leanok
Import the existing declaration.
\end{proof}

\begin{theorem}[GMT: finite restrictions are regular]
\label{gmt:finite-restriction}
\lean{BorelRegularOuterMeasure.restrict_isRadon}
\leanok
On a locally compact, Hausdorff, second-countable Borel space, $\mu E<\infty$ implies that $\mu|E$ is regular. The source does not require $E$ to be Borel.
\end{theorem}
\begin{proof}[Formalized in FormalizingGMT]
\leanok
Import the existing declaration.
\end{proof}

\begin{definition}[GMT: Hausdorff content]
\label{gmt:hausdorff-content}
\lean{hausdorffContent,hausdorffContentInfty}
\leanok
Hausdorff content is the infimum of countable covering diameter sums, with or without a diameter cutoff. For positive exponents, empty cover elements have zero contribution; the cutoff version uses an explicit nonempty guard.
\end{definition}

\begin{lemma}[GMT: content comparison and stabilization]
\label{gmt:content-comparison}
\lean{hausdorffContentInfty_le_hausdorffContent,hausdorffContent_le_hausdorffMeasure,hausdorffContent_eq_hausdorffContentInfty_of_ediam_le}
\leanok
\uses{gmt:hausdorff-content}
Unrestricted content is bounded above by cutoff content, which is bounded above by Hausdorff measure. For positive exponent and a positive cutoff at least the diameter of a set, cutoff content equals unrestricted content.
\end{lemma}
\begin{proof}[Formalized in FormalizingGMT]
\leanok
Import the existing declaration.
\end{proof}

\begin{theorem}[GMT: finite Hausdorff restrictions]
\label{gmt:hausdorff-regular}
\lean{HausdorffRestrict.toRadonOuterMeasure}
\leanok
\uses{gmt:hausdorff-content}
On a locally compact, second-countable Hausdorff metric space, restriction of $\mathcal H^s$ to a set of finite $\mathcal H^s$ measure is regular; use the source's Caratheodory-measurability and $s\geq0$ interface.
\end{theorem}
\begin{proof}[Formalized in FormalizingGMT]
\leanok
\uses{gmt:finite-restriction}
Import the existing declaration.
\end{proof}

\begin{theorem}[GMT: positive upper content density inside]
\label{gmt:content-density}
\lean{hausdorffContentInfty_upperDensity_ge_ae_mem}
\leanok
\uses{def:density-api,gmt:hausdorff-content}
For $s\geq0$ and $\mathcal H^s(E)<\infty$ in a metric space, the upper diameter-normalized density of restricted Hausdorff content is at least $2^{-s}$ at $\mathcal H^s$-almost every point of $E$.
\end{theorem}
\begin{proof}[Formalized in FormalizingGMT]
\leanok
\uses{gmt:content-comparison}
Import the existing declaration.
\end{proof}

\begin{theorem}[GMT: positive upper Hausdorff density inside]
\label{gmt:hausdorff-density-lower}
\lean{hausdorffMeasure_upperDensity_ge_ae_mem}
\leanok
\uses{def:density-api}
For $s\geq0$ and $\mathcal H^s(E)<\infty$, $\Theta^{*s}(\mathcal H^s|E,x)\geq2^{-s}$ at $\mathcal H^s$-almost every $x\in E$. This is a lower bound on the UPPER density, not on the lower density.
\end{theorem}
\begin{proof}[Formalized in FormalizingGMT]
\leanok
\uses{gmt:content-density,gmt:content-comparison}
Import the existing declaration.
\end{proof}

\begin{theorem}[GMT: upper Hausdorff density at most one]
\label{gmt:hausdorff-density-upper}
\lean{hausdorffMeasure_upperDensity_le_one_ae_mem}
\leanok
\uses{def:density-api}
If $E$ is Caratheodory measurable with finite $\mathcal H^s$ measure, $s\geq0$, and the metric space is locally compact and second countable, then $\Theta^{*s}(\mathcal H^s|E,x)\leq1$ for $\mathcal H^s$-almost every $x\in E$.
\end{theorem}
\begin{proof}[Formalized in FormalizingGMT]
\leanok
\uses{gmt:hausdorff-regular}
Import the existing declaration.
\end{proof}

\begin{theorem}[GMT: zero upper density outside]
\label{gmt:hausdorff-density-outside}
\lean{hausdorffMeasure_upperDensity_eq_zero_ae_notMem}
\leanok
\uses{def:density-api}
Under the same local compactness, second countability, measurability and finite-measure hypotheses, the upper $s$-density of $\mathcal H^s|E$ is zero for $\mathcal H^s$-almost every $x\notin E$.
\end{theorem}
\begin{proof}[Formalized in FormalizingGMT]
\leanok
\uses{gmt:hausdorff-regular}
Import the existing declaration.
\end{proof}

\begin{definition}[GMT: vague convergence]
\label{gmt:vague}
\lean{MeasureTheory.Measure.WeaklyConverges}
\leanok
Convergence of a sequence of measures is tested against every real continuous compactly supported function. This does not require finite total mass.
\end{definition}

\begin{definition}[GMT: blow-up map]
\label{gmt:blowup}
\lean{blowUpMap}
\leanok
The blow-up map on Euclidean space is $T_{a,r}(z)=r^{-1}(z-a)$. Geometric use assumes $r>0$.
\end{definition}

\begin{lemma}[GMT: measurable blow-ups and ball identities]
\label{gmt:blowup-balls}
\lean{measurable_blowUpMap,map_blowUp_apply_ball,map_blowUp_apply_closedBall}
\leanok
\uses{gmt:blowup}
Blow-ups are measurable. For $r>0$, the pushforward mass of $B(0,\rho)$ or $\overline B(0,\rho)$ equals the original mass at centre $a$ and radius $r\rho$.
\end{lemma}
\begin{proof}[Formalized in FormalizingGMT]
\leanok
Import the existing declaration.
\end{proof}

\begin{lemma}[GMT: regularity of scaled blow-ups]
\label{gmt:blowup-regular}
\lean{regular_map_blowUp,regular_smul_map_blowUp}
\leanok
\uses{gmt:blowup}
For $r\ne0$, blow-up pushforwards preserve regularity; multiplying by a finite extended nonnegative scalar preserves it too.
\end{lemma}
\begin{proof}[Formalized in FormalizingGMT]
\leanok
Import the existing declaration.
\end{proof}

\begin{definition}[GMT: tangent measures]
\label{def:tangent-api}
\lean{IsTangentMeasure}
\leanok
\uses{gmt:vague,gmt:blowup}
A tangent is a nonzero regular measure obtained as a vague limit of $c_i(T_{a,r_i})_\#\mu$ with $r_i>0$, $r_i\to0$, and $0<c_i<\infty$. The scalar sequence is arbitrary.
\end{definition}

\begin{theorem}[GMT: compact/open portmanteau]
\label{gmt:portmanteau}
\lean{MeasureTheory.Measure.weaklyConverges_iff_compactOpenBounds}
\leanok
\uses{gmt:vague}
For regular Euclidean measures, vague convergence is equivalent to the upper bounds on compact sets and lower bounds on open sets.
\end{theorem}
\begin{proof}[Formalized in FormalizingGMT]
\leanok
Import the existing declaration.
\end{proof}

\begin{theorem}[GMT: convergence on bounded continuity sets]
\label{gmt:continuity-sets}
\lean{MeasureTheory.Measure.weaklyConverges_iff_boundedContinuitySets}
\leanok
\uses{gmt:vague}
For regular Euclidean measures, vague convergence is equivalent to convergence of masses on bounded measurable sets whose boundary has zero limit mass.
\end{theorem}
\begin{proof}[Formalized in FormalizingGMT]
\leanok
\uses{gmt:portmanteau}
Import the existing declaration.
\end{proof}

\begin{theorem}[GMT: compact mass bounds give a subsequence]
\label{thm:compactness}
\lean{MeasureTheory.exists_vaguelyConvergent_subseq_of_compact_bounded}
\leanok
\uses{gmt:vague}
Uniform mass bounds on every compact set give a strictly indexed subsequence converging vaguely to a regular measure. Nonzeroness of the limit must be proved separately.
\end{theorem}
\begin{proof}[Formalized in FormalizingGMT]
\leanok
Import the existing declaration.
\end{proof}

\begin{lemma}[GMT: density bounds imply support and doubling]
\label{gmt:doubling-bounds}
\lean{mem_support_of_ball_lower_bound,limsup_ball_ratio_lt_top_of_ball_bounds_at}
\leanok
\uses{def:density-api}
At a point with positive lower and finite upper power bounds on sufficiently small closed balls, the point belongs to the support and its upper small-scale open-ball doubling ratio is finite.
\end{lemma}
\begin{proof}[Formalized in FormalizingGMT]
\leanok
Import the existing declaration.
\end{proof}

\begin{theorem}[GMT: origin lies in tangent support]
\label{gmt:origin-support}
\lean{zero_mem_support_of_isTangentMeasure}
\leanok
\uses{def:tangent-api}
For a regular measure at a support point with finite upper doubling ratio, every tangent has zero in its support.
\end{theorem}
\begin{proof}[Formalized in FormalizingGMT]
\leanok
\uses{gmt:portmanteau,gmt:blowup-balls}
Import the existing declaration.
\end{proof}

\begin{theorem}[GMT: mass-normalized tangent sequence]
\label{gmt:normalized-tangent}
\lean{exists_normalized_blowUp_weaklyConverges_of_isTangentMeasure}
\leanok
\uses{def:tangent-api}
Under the same support and doubling assumptions, each tangent is the vague limit of a fixed positive finite multiple of the ball-mass-normalized blow-ups along strictly decreasing radii.
\end{theorem}
\begin{proof}[Formalized in FormalizingGMT]
\leanok
\uses{gmt:origin-support,gmt:continuity-sets}
Import the existing declaration.
\end{proof}

\begin{theorem}[GMT: tangent existence under doubling]
\label{thm:doubling}
\lean{exists_subseq_blowUp_weaklyConverges_tangentMeasure}
\leanok
\uses{def:tangent-api}
At a support point of a regular measure with finite upper doubling ratio, every positive sequence of radii tending to zero has a subsequence of ball-mass-normalized blow-ups converging to a nonzero tangent (Mattila 14.3).
\end{theorem}
\begin{proof}[Formalized in FormalizingGMT]
\leanok
\uses{thm:compactness,gmt:blowup-balls,gmt:blowup-regular,gmt:portmanteau}
Import the existing declaration.
\end{proof}

\begin{theorem}[GMT: restriction preserves tangents]
\label{gmt:restriction-tangents}
\lean{isTangentMeasure_restrict_iff}
\leanok
\uses{def:tangent-api}
For Borel $B$, regular $\mu$ and $\mu|B$, and a support point $a$ with $\mu(B(a,r)\setminus B)/\mu(B(a,r))\to0$, restriction to $B$ preserves the tangent set at $a$ (Mattila 14.5).
\end{theorem}
\begin{proof}[Formalized in FormalizingGMT]
\leanok
\uses{gmt:blowup-regular}
Import the existing declaration.
\end{proof}

\begin{theorem}[GMT: density change preserves tangents]
\label{thm:matching}
\lean{isTangentMeasure_iff_ae_of_density}
\leanok
\uses{def:tangent-api}
If regular measures satisfy $\lambda(A)=\int_A\varphi\,d\mu$ on Borel sets for a nonnegative real locally integrable $\varphi$, their tangent sets coincide at $\lambda$-almost every point (Mattila 14.6).
\end{theorem}
\begin{proof}[Formalized in FormalizingGMT]
\leanok
\uses{gmt:blowup-regular}
Import the existing declaration.
\end{proof}

\begin{lemma}[GMT: closed good sets for density bounds]
\label{gmt:goodsets}
\lean{positiveFiniteDensitySet,RatioOfDensities,goodSet,isClosed_goodSet,exists_goodSet_mem}
\leanok
\uses{def:density-api}
Positive finite density-bound points admit a countable stratification by closed good sets with uniform small-radius power bounds. The rational bounds may approximate the lower/upper density ratio.
\end{lemma}
\begin{proof}[Formalized in FormalizingGMT]
\leanok
\uses{gmt:density-bounds}
Import the existing declaration.
\end{proof}

\begin{lemma}[GMT: one null set of bad density points]
\label{gmt:density-points}
\lean{exists_null_of_not_density_point}
\leanok
For a Borel set $B$ and regular Euclidean $\mu$, one null exceptional set can be removed so that every remaining $a\in B$ satisfies, for each $\gamma>0$, $\mu(\overline B(a,r)\setminus B)\leq\gamma\mu(\overline B(a,r))$ eventually as $r\downarrow0$.
\end{lemma}
\begin{proof}[Formalized in FormalizingGMT]
\leanok
Import the existing declaration.
\end{proof}

\begin{theorem}[GMT: transfer bounds at a density point]
\label{gmt:transfer-ball-bounds}
\lean{tangent_ball_bounds_of_density_point}
\leanok
\uses{def:tangent-api}
Uniform small-ball bounds $td\,r^s\leq\mu(\overline B(y,r))\leq d\,r^s$ on a set transfer at its density points to bounds $tc\,r^s\leq\nu(\overline B(z,r))\leq c\,r^s$ at all support centres and all positive radii of every tangent; $c$ is positive and finite.
\end{theorem}
\begin{proof}[Formalized in FormalizingGMT]
\leanok
\uses{gmt:portmanteau,gmt:blowup-balls,gmt:blowup-regular}
Import the existing declaration.
\end{proof}

\begin{theorem}[GMT: tangent bounds from density ratio]
\label{gmt:mattila-14-7-1}
\lean{mattila_14_7_1}
\leanok
\uses{def:tangent-api,def:density-api}
Outside one $\mu$-null set, every tangent at a positive finite density-bound point has support-centred power bounds with ratio $\Theta_*^s/\Theta^{*s}$ (Mattila 14.7(1)).
\end{theorem}
\begin{proof}[Formalized in FormalizingGMT]
\leanok
\uses{gmt:goodsets,gmt:density-points,gmt:transfer-ball-bounds}
Import the existing declaration.
\end{proof}

\begin{theorem}[GMT: uniform support bounds and all tangents]
\label{gmt:mattila-14-7-4}
\lean{mattila_14_7_4}
\leanok
\uses{def:tangent-api}
If one pair of positive power bounds holds at every support point below one fixed radius, then at every support point tangents exist and all of them satisfy the corresponding global support-centred power bounds (Mattila 14.7(4)).
\end{theorem}
\begin{proof}[Formalized in FormalizingGMT]
\leanok
\uses{thm:doubling,gmt:portmanteau,gmt:blowup-balls}
Import the existing declaration.
\end{proof}

\begin{definition}[GMT: uniform measures]
\label{def:uniform-api}
\lean{IsSUniform}
\leanok
A nonzero regular measure is $s$-uniform when $s>0$ and $\nu(\overline B(y,r))=cr^s$ for one $c>0$, all $y\in\spt\nu$ and all $r>0$.
\end{definition}

\begin{definition}[GMT: positive finite existing density set]
\label{gmt:density-exists-set}
\lean{positiveFiniteDensityExistsSet}
\leanok
\uses{def:density-api}
This set is defined using \texttt{HasDensity} together with a positive finite value of the density. It requires an adapter from our liminf/limsup formulation.
\end{definition}

\begin{lemma}[GMT: existing density gives ratio one and bounds]
\label{gmt:uniform-density-adapter}
\lean{SUniformAux.mem_positiveFiniteDensitySet,SUniformAux.ball_bounds_of_mem}
\leanok
\uses{gmt:density-exists-set,gmt:goodsets}
Membership in the positive-finite-existing-density set implies membership in the two-sided density-bound set with ratio one. Such points have positive lower and finite upper small-ball power bounds.
\end{lemma}
\begin{proof}[Formalized in FormalizingGMT]
\leanok
\uses{gmt:density-bounds}
Import the existing declaration.
\end{proof}

\begin{theorem}[GMT: all tangents are uniform]
\label{thm:uniform-upstream}
\lean{mattila_14_9}
\leanok
\uses{def:uniform-api,def:tangent-api,gmt:density-exists-set}
For $s>0$ and regular $\mu$, outside one null set every tangent at every positive finite existing-density point is $s$-uniform and has the origin in its support (Mattila 14.9).
\end{theorem}
\begin{proof}[Formalized in FormalizingGMT]
\leanok
\uses{gmt:mattila-14-7-1,gmt:uniform-density-adapter,gmt:doubling-bounds,gmt:origin-support}
Import the existing declaration.
\end{proof}

\begin{lemma}[GMT: tangents of uniform measures]
\label{gmt:uniform-tangents}
\lean{SUniformAux.exists_isTangentMeasure,SUniformAux.isSUniform_of_isTangentMeasure}
\leanok
\uses{def:uniform-api,def:tangent-api}
At every support point of an $s$-uniform measure tangents exist, and every such tangent is $s$-uniform with zero in its support.
\end{lemma}
\begin{proof}[Formalized in FormalizingGMT]
\leanok
\uses{gmt:mattila-14-7-4,gmt:origin-support}
Import the existing declaration.
\end{proof}

\begin{lemma}[GMT: uniform measure in a half-space]
\label{gmt:uniform-halfspace}
\lean{SUniformAux.exists_halfSpace}
\leanok
\uses{def:uniform-api}
A uniform measure whose support is not the whole ambient space produces a uniform measure supported in a half-space and containing zero in its support. The proof constructs a tangent, but the exported conclusion does not retain that tangent witness.
\end{lemma}
\begin{proof}[Formalized in FormalizingGMT]
\leanok
\uses{gmt:uniform-tangents,gmt:portmanteau}
Import the existing declaration.
\end{proof}

\begin{lemma}[GMT: bounded radial integrals agree]
\label{gmt:uniform-radial}
\lean{SUniformAux.measure_closedBall_eq_of_mem_support,SUniformAux.setIntegral_comp_dist_eq}
\leanok
\uses{def:uniform-api}
For an $s$-uniform measure, same-radius balls centred on the support have equal mass. The integrals of any continuous radial function over those closed balls agree. Passing to global nonnegative Borel radial tests and identifying the Euclidean radial distribution still need adapters.
\end{lemma}
\begin{proof}[Formalized in FormalizingGMT]
\leanok
Import the existing declaration.
\end{proof}

\begin{lemma}[GMT: half-space to hyperplane]
\label{gmt:uniform-hyperplane}
\lean{SUniformAux.exists_hyperplane}
\leanok
\uses{def:uniform-api}
An $s$-uniform measure containing zero in its support and supported in a half-space produces an $s$-uniform measure supported in a hyperplane. The result does not export the lineage of the blow-ups used in its proof.
\end{lemma}
\begin{proof}[Formalized in FormalizingGMT]
\leanok
\uses{gmt:uniform-tangents,gmt:uniform-radial}
Import the existing declaration.
\end{proof}

\begin{lemma}[GMT: transport from a hyperplane]
\label{gmt:hyperplane-transport}
\lean{SUniformAux.exists_isSUniform_of_hyperplane}
\leanok
\uses{def:uniform-api}
An $s$-uniform measure on $\R^{n+1}$ supported in a linear hyperplane yields an $s$-uniform measure on $\R^n$. The proof uses a linear isometry of the hyperplane; the existential interface does not expose it.
\end{lemma}
\begin{proof}[Formalized in FormalizingGMT]
\leanok
Import the existing declaration.
\end{proof}

\begin{lemma}[GMT: lower-dimensional uniform measure]
\label{gmt:uniform-lower-dim}
\lean{SUniformAux.exists_lower_dim}
\leanok
\uses{def:uniform-api}
A uniform measure on $\R^{n+1}$ whose support is not the whole space yields a uniform measure with the same exponent on $\R^n$. This is not yet the moving-centre pseudotangent theorem needed for Part A.
\end{lemma}
\begin{proof}[Formalized in FormalizingGMT]
\leanok
\uses{gmt:uniform-halfspace,gmt:uniform-hyperplane,gmt:hyperplane-transport}
Import the existing declaration.
\end{proof}

\begin{theorem}[GMT: full-dimensional uniform measures are flat]
\label{gmt:uniform-full-dimension}
\lean{mattila_14_11}
\leanok
\uses{def:uniform-api}
An $n$-uniform measure on $\R^n$ is a positive finite scalar multiple of Lebesgue measure (Mattila 14.11). The source definition of uniformity implies $n>0$.
\end{theorem}
\begin{proof}[Formalized in FormalizingGMT]
\leanok
\uses{gmt:uniform-lower-dim}
Import the existing declaration.
\end{proof}

\section{Additional reuse exposed by the proof audit}

\begin{lemma}[GMT: lower growth excludes full support]
\label{gmt:support-not-full}
\lean{support_ne_univ_of_lower_growth}
\leanok
A regular measure on $\R^n$ with a common positive support-centred lower $s$-power bound has non-full support if $s<n$.
\end{lemma}
\begin{proof}[Existing formalization]
\leanok
\uses{def:uniform-api}
Import the linked declaration at the pinned source revision.
\end{proof}

\begin{lemma}[GMT: quadratic barycentre estimate]
\label{gmt:barycenter-quadratic-bound}
\lean{SUniformAux.abs_inner_integral_le}
\leanok
For a uniform measure with zero in its support and each positive radius $r$, the inner product of a support point $y$ with the vector integral over $\overline B(0,r)$ is bounded in absolute value by $C|y|^2$ for $|y|\leq r$.
\end{lemma}
\begin{proof}[Existing formalization]
\leanok
\uses{def:uniform-api,gmt:uniform-radial}
Import the linked declaration at the pinned source revision.
\end{proof}

\begin{lemma}[GMT: tangent support from a quadratic estimate]
\label{gmt:quadratic-support-tangent}
\lean{SUniformAux.support_tangent_subset_of_estimate}
\leanok
If $|\langle y,w\rangle|\leq C|y|^2$ holds on the support near zero, every tangent at zero is supported in $w^\perp$. The estimate has nonnegative finite $C$ and a positive radius of validity.
\end{lemma}
\begin{proof}[Existing formalization]
\leanok
\uses{def:uniform-api,def:tangent-api}
Import the linked declaration at the pinned source revision.
\end{proof}

\begin{lemma}[GMT: vanishing half-space integrals force a hyperplane]
\label{gmt:zero-barycenter-hyperplane}
\lean{SUniformAux.support_subset_of_integral_inner_eq_zero}
\leanok
For a uniform measure supported where $\langle z,e\rangle\geq0$, vanishing of the integral of this inner product over every positive-radius centred closed ball forces its support into $e^\perp$.
\end{lemma}
\begin{proof}[Existing formalization]
\leanok
\uses{def:uniform-api}
Import the linked declaration at the pinned source revision.
\end{proof}

\begin{lemma}[GMT: compact-test estimates and perturbation]
\label{gmt:compact-test-perturbation}
\lean{norm_integral_cc_le,weaklyConverges_of_tendsto_integral_sub}
\leanok
A compact-test integral is bounded by the test supremum norm times the finite mass of a containing support set. Two sequences whose compact-test integral differences tend to zero have the same vague limits.
\end{lemma}
\begin{proof}[Existing formalization]
\leanok
\uses{gmt:vague}
Import the linked declaration at the pinned source revision.
\end{proof}

\begin{lemma}[GMT: negligible discarded mass in normalized blow-ups]
\label{gmt:scaled-discarded-mass}
\lean{tendsto_nhdsGT_zero_of_pos,eventually_smul_measure_closedBall_le,tendsto_mul_measure_diff_zero}
\leanok
Positive radii tending to zero tend through the positive neighborhood filter. Convergent scaled blow-ups have eventual local ball bounds. At a support point where a restriction has relative density one, those bounds make the scaled discarded ball mass tend to zero.
\end{lemma}
\begin{proof}[Existing formalization]
\leanok
\uses{gmt:blowup-balls,gmt:portmanteau}
Import the linked declaration at the pinned source revision.
\end{proof}

\begin{lemma}[mathlib: measurable hull intersection identity]
\label{mathlib:measurable-hull}
\lean{MeasureTheory.Measure.measure_toMeasurable_inter}
\leanok
For a set $D$ of finite measure and measurable $U$, $\mu(\operatorname{toMeasurable}_\mu D\cap U)=\mu(D\cap U)$. No measurability of $D$ is required.
\end{lemma}
\begin{proof}[Existing formalization]
\leanok
Import the linked declaration at the pinned source revision.
\end{proof}

\begin{lemma}[mathlib: disjoint almost-everywhere ball covering]
\label{mathlib:besicovitch-cover}
\lean{Besicovitch.exists_disjoint_closedBall_covering_ae}
\leanok
For an s-finite measure in a space with the Besicovitch covering property, a fine family of positive-radius closed balls admits a countable disjoint subfamily covering the target set almost everywhere, with prescribed positive upper radius bounds.
\end{lemma}
\begin{proof}[Existing formalization]
\leanok
Import the linked declaration at the pinned source revision.
\end{proof}

\begin{lemma}[mathlib: differentiation of measures and restrictions]
\label{mathlib:besicovitch-differentiation}
\lean{Besicovitch.ae_tendsto_rnDeriv,Besicovitch.ae_tendsto_measure_inter_div_of_measurableSet,Besicovitch.ae_tendsto_measure_inter_div}
\leanok
For locally finite measures on finite-dimensional real normed spaces, closed-ball measure ratios converge almost everywhere to the Radon--Nikodym density. Restrictions have relative density one almost everywhere on the restricted set; the arbitrary-set version uses the restricted measure.
\end{lemma}
\begin{proof}[Existing formalization]
\leanok
Import the linked declaration at the pinned source revision.
\end{proof}

\begin{lemma}[mathlib: disjoint selection with enlarged ball covering]
\label{mathlib:vitali-enlargement}
\lean{Vitali.exists_disjoint_subfamily_covering_enlargement_ball}
\leanok
A bounded-radius ball family admits a disjoint subfamily such that each original ball is contained in an enlargement of a selected ball by any fixed factor greater than three. This is distinct from a disjoint almost-everywhere covering without enlargement.
\end{lemma}
\begin{proof}[Existing formalization]
\leanok
Import the linked declaration at the pinned source revision.
\end{proof}



\section{Further generic mathlib infrastructure}
\begin{lemma}[Derivative-controlled closed covers and measurable partitions]
\label{mathlib:derivative-partitions}
\lean{exists_closed_cover_approximatesLinearOn_of_hasFDerivWithinAt, exists_partition_approximatesLinearOn_of_hasFDerivWithinAt}
\leanok
For a differentiable map between finite-dimensional normed spaces, mathlib supplies
countable pieces with uniform linear approximation by actual derivatives.
These generic results apply in codimension; the same-dimensional change-of-variables
formula in that file does not itself supply the graph area formula.
\end{lemma}
\begin{proof}[Existing upstream proof]
\leanok
Source: \texttt{Mathlib/MeasureTheory/Function/Jacobian.lean}.
\end{proof}

\begin{definition}[Intrinsic standard Gaussian measure]
\label{mathlib:standard-gaussian}
\lean{ProbabilityTheory.stdGaussian, ProbabilityTheory.stdGaussian_map, ProbabilityTheory.stdGaussian_eq_map_pi_orthonormalBasis}
\leanok
Mathlib defines the standard Gaussian measure on a finite-dimensional real inner
product space, proves its behavior under linear isometries, and identifies it with
the orthonormal-coordinate product Gaussian. Use this measure on a plane subtype
instead of defining a second Gaussian distribution.
\end{definition}

\begin{lemma}[Uniform measures cannot exceed ambient dimension]
\label{gmt:uniform-dimension-bound}
\lean{SUniformAux.not_isSUniform_of_lt}
\leanok
\uses{def:uniform-api}
There is no $s$-uniform measure on $\mathbb R^n$ when $n<s$.
\end{lemma}
\begin{proof}[Existing upstream proof]
\leanok
FormalizingGMT supplies the obstruction directly. Together with tangent existence
and almost-everywhere uniformity, it proves the excessive-dimension endpoint.
\end{proof}

\begin{lemma}[GMT: closed sets controlling all balls through a point]
\label{gmt:goodballs}
\lean{goodBallSet, isClosed_goodBallSet}
\leanok
\uses{def:density-api}
The set of points for which every sufficiently small closed ball containing
the point satisfies a fixed upper power bound is closed.
\end{lemma}
\begin{proof}[Existing upstream proof]
\leanok
Enlarge a ball slightly, approximate the centre point by good points,
and let the enlargement tend to zero.
\end{proof}

\begin{lemma}[mathlib: mass-distribution bound for Hausdorff measure]
\label{mathlib:hausdorff-mass-bound}
\lean{MeasureTheory.Measure.le_hausdorffMeasure}
\leanok
A measure bounded by the diameter to a real power on every sufficiently
small set is at most the corresponding Hausdorff measure.
\end{lemma}
\begin{proof}[Existing mathlib proof]
\leanok
Use the maximal outer measure defining the metric Hausdorff construction.
\end{proof}

\begin{lemma}[mathlib: compact tests determine regular measures]
\label{mathlib:regular-measure-ext}
\lean{MeasureTheory.Measure.ext_of_integral_eq_on_compactlySupported}
\leanok
On a locally compact Hausdorff space with its Borel measurable structure,
two regular measures are equal if their integrals agree against every
real-valued compactly supported continuous function.
\end{lemma}
\begin{proof}[Existing mathlib proof]
\leanok
Use regularity and continuous compactly supported approximations to compact sets.
\end{proof}


\section{Other reuse and limits of the audit}
Mathlib already supplies Hausdorff measure, Lipschitz distortion of Hausdorff
measure, Besicovitch differentiation, Radon--Nikodym derivatives,
\texttt{LipschitzWith.ae\_differentiableAt}, orthogonal projections,
Gaussian integrals, and Stone--Weierstrass tools. Source paths and exact
candidate declarations are recorded in \texttt{reuse.json}.
No rectifiability definition was found in the pinned mathlib/GMT search.
The algebraic \texttt{Module.Grassmannian} is a quotient-rank construction;
it does not supply the compact metric space of Euclidean planes needed here.
Do not assume that this topological API is already available.

FormalizingGMT's Marstrand dimension reduction contains useful hyperplane
and lower-dimensional uniform-measure lemmas. These are candidates for reuse
inside the pseudotangent construction, but their conclusions alone do not
record a moving-centre blow-up sequence. That missing compatibility is a
proof obligation. The integer-dimension theorem itself is optional because
$m$ is already a natural number in our target.

\chapter{Reusable definitions and normalization}

\begin{definition}[Countable Lipschitz coverage]
\label{def:coverage}
\lean{Preiss.CountablyLipschitzCovered}
A set $S$ is countably Lipschitz covered in dimension $m$ if a countable
family of Lipschitz maps $f:\R^m\to X$ covers $S$. The family may be empty.
No measurability or finite-measure hypothesis is part of this predicate.
\end{definition}

\begin{definition}[Rectifiable sets and measures]
\label{def:rectifiable}
\lean{Preiss.CountablyRectifiable,Preiss.IsRectifiableMeasure}
\uses{def:coverage}
A set $S$ is countably $m$-rectifiable if $\Hm(S\setminus T)=0$ for some
countably Lipschitz covered $T$. A measure $\mu$ is $m$-rectifiable if
$\mu\ll\Hm$ and $\mu$ is concentrated on a countably Lipschitz covered set.
Merely being carried by Lipschitz images is not enough: it would admit atoms
in positive dimension. Use mathlib's Hausdorff measure without redefining it.
\end{definition}

\begin{lemma}[Carrier calculus]
\label{lem:carrier-calculus}
\uses{def:coverage,def:rectifiable}
Coverage is hereditary and closed under countable unions and Lipschitz
composition. In a Euclidean target a countable Lipschitz image union is a
countable union of compact images of closed balls, hence Borel and
$\Hm$-sigma-finite. Hausdorff-null remainders preserve rectifiability.
\end{lemma}
\begin{proof}[Planned proof]
Use compact exhaustions of $\R^m$ and the existing Lipschitz Hausdorff
measure bound. This is an early reusable API milestone.
\end{proof}

\begin{lemma}[Lipschitz images versus graphs]
\label{lem:graphs}
\uses{def:rectifiable}
In Euclidean space the image-based definition agrees, modulo
$\Hm$-null sets, with countable coverage by $m$-dimensional Lipschitz graphs,
as in De Lellis Definition 4.1.
\end{lemma}
\begin{proof}[Planned proof]
Decompose the full-rank differentiability locus into graph pieces and show
the rank-deficient image is null. Reuse Rademacher; isolate the null-image
lemma (De Lellis 4.11). The full area formula and the reverse implication
of Theorem 4.8 are useful extensions, not mandatory inputs to Preiss.
Coordinatewise Lipschitz extension suffices; an optimal Kirszbraun constant
is unnecessary. This equivalence must be proved before claiming literal
agreement with the source's graph convention.
\end{proof}

\begin{lemma}[Density representation]
\label{lem:representation}
\lean{PreissBlueprint.rectifiable_iff_density_representation}
\uses{def:rectifiable}
A locally finite rectifiable measure can be written
$f\,(\Hm|S)$,
where $S$ is Borel and rectifiable, $\Hm|S$ is sigma-finite, and
$f:X\to[0,\infty]$ is measurable; conversely such a measure is rectifiable.
\end{lemma}
\begin{proof}[Planned proof]
\uses{lem:carrier-calculus}
Apply Radon--Nikodym to the restricted measure on the carrier.
Ambient $\Hm$ on $\R^n$ is not sigma-finite for $m<n$; applying the theorem
to ambient Hausdorff measure would be a mistake. For locally finite $\mu$,
the density can additionally be chosen finite almost everywhere on the carrier.
\end{proof}

\begin{definition}[Positive finite density]
\label{def:positive-density}
\lean{Preiss.HasPositiveFiniteDensityAt,Preiss.HasPositiveFiniteDensitiesAt}
\uses{def:density-api}
The first predicate means $0<\Theta_*^m=\Theta^{*m}<\infty$ in the
upstream convention. The second means only
$0<\Theta_*^m\leq\Theta^{*m}<\infty$.
\end{definition}

\begin{lemma}[Normalization bridges]
\label{lem:density-bridge}
\lean{PreissBlueprint.density_iff_upstream,PreissBlueprint.density_iff_openBall_limit}
\uses{def:positive-density}
Positive finite equal densities are equivalent to membership in the existing
positive-finite-density-exists set, and to a positive finite limit of
$\mu(B(x,r))/r^m$.
\end{lemma}
\begin{proof}[Planned proof]
\uses{gmt:density-limit,gmt:density-bounds,gmt:density-exists-set}
Use the complete order topology of ENNReal. Squeeze open and closed balls
using radii $(1\pm\varepsilon)r$, then let $\varepsilon\downarrow0$.
The radius-normalized limit is $2^m$ times the diameter-normalized one.
De Lellis further divides by $\omega_m$; the resulting constant is positive
and finite. Equality, positivity, and finiteness are invariant under this
change, but exact constants in uniform-measure lemmas must be transported.
\end{proof}

\begin{definition}[Planes, flat measures, and dimensional scaling]
\label{def:flat}
\lean{Preiss.Plane,Preiss.IsFlat,Preiss.dimensionalBlowup,Preiss.IsDimensionalTangent,Preiss.IsTangentAtInfinity,Preiss.IsUnitUniform,Preiss.IsConical}
\uses{def:tangent-api,def:uniform-api}
An $m$-plane is a real submodule of dimension $m$. A flat measure is
$c\Hm|V$ for such a plane and $0<c<\infty$. Put
$D^m_{x,r}\mu=r^{-m}(T_{x,r})_\#\mu$ for $r>0$.
Dimensional tangents use these fixed factors at radii tending to zero;
tangents at infinity use radii tending to infinity and centre zero.
A unit uniform measure is an existing $m$-uniform measure with
$0\in\spt\nu$ and $\nu(\overline B(0,1))=1$.
A conical measure satisfies $D^m_{0,r}\nu=\nu$ for every $r>0$.
\end{definition}

\begin{lemma}[Plane and scaling API]
\label{lem:plane-scaling}
\uses{def:flat}
Flat measures are nonzero regular measures; for $m>0$ they are uniform.
Scaling, recentering at a support point, and linear isometries preserve the
appropriate classes. Hausdorff restriction to a plane is a fixed positive
multiple of its induced Lebesgue measure. Unit normalization removes that
constant. Blow-up composition obeys
$D^m_{y,s}(D^m_{x,r}\mu)=D^m_{x+ry,rs}\mu$.
\end{lemma}
\begin{proof}[Planned proof]
\uses{gmt:blowup-balls,gmt:blowup-regular}
Reuse the pushforward ball identities and regularity lemmas; establish the plane Hausdorff normalization and composition adapters.
\end{proof}

\begin{lemma}[Comparison of tangent conventions]
\label{lem:tangent-ray}
\lean{PreissBlueprint.tangent_iff_dimensional_ray}
\uses{def:positive-density,def:flat}
At a positive finite density point of a locally finite measure, an arbitrary
scalar tangent is exactly a positive finite scalar multiple of a dimensional
tangent. These two sets of tangents are not equal.
\end{lemma}
\begin{proof}[Planned proof]
\uses{lem:density-bridge,thm:compactness,lem:plane-scaling,gmt:normalized-tangent,gmt:continuity-sets}
Control the scalar factors by masses of continuity balls, extract their
positive finite limit, and pass scaling through compactly supported tests.
Establish uniform local bounds before passing to the limit.
\end{proof}

\chapter{Rectifiability from flat tangents}

\begin{lemma}[Absolute continuity from upper density]
\label{lem:absolute-continuity}
\lean{PreissBlueprint.absolutelyContinuous_hausdorff_of_upper_density_finite}
\uses{def:density-api}
If $\mu$ is locally finite and $\Theta^{*m}(\mu,x)<\infty$ at
$\mu$-almost every point, then $\mu\ll\Hm$.
\end{lemma}
\begin{proof}[Planned proof]
\uses{gmt:vitali}
Stratify by uniform small-ball upper bounds and use Hausdorff covers.
Reuse covering tools; this direction needs the upper bound, not a lower bound.
\end{proof}

\begin{lemma}[Finite Hausdorff pieces]
\label{lem:finite-pieces}
\uses{def:positive-density}
Under positive finite lower and upper density bounds almost everywhere,
there are countably many Borel pieces carrying $\mu$, each of finite
$\Hm$ measure, on which $\mu$ and $\Hm$ are mutually absolutely continuous
with density bounded above and below by positive constants.
\end{lemma}
\begin{proof}[Planned proof]
\uses{lem:absolute-continuity,gmt:goodsets,gmt:finite-restriction,gmt:hausdorff-regular,gmt:vitali}
Stratify by radii, locations and two density bounds; use Vitali comparisons.
Derive the restrictions' regularity from finite measure. Do not assert
regularity of arbitrary infinite restrictions without justification.
\end{proof}

\begin{lemma}[Cone-to-graph lemma]
\label{lem:cone-graph}
\uses{def:coverage}
If a set $F$ satisfies $\|Q_V(y-x)\|\leq L\|P_V(y-x)\|$ for all
$x,y\in F$, for one $m$-plane $V$ and finite $L\geq0$, then $F$ lies
in a Lipschitz graph over $V$. This is De Lellis Lemma 4.7.
\end{lemma}
\begin{proof}[Planned proof]
Projection onto $V$ is injective on $F$; its inverse defines a Lipschitz
graph on the projected set. Extend coordinatewise, then identify $V$ with $\R^m$.
\end{proof}

\begin{lemma}[Approximate cone criterion]
\label{lem:cone-criterion}
\uses{def:rectifiable}
For a Borel set of finite $\Hm$ measure, positive lower density and
vanishing $r^{-m}$-normalized mass outside a fixed cone at almost every point
imply countable rectifiability (De Lellis Proposition 4.6).
\end{lemma}
\begin{proof}[Planned proof]
\uses{lem:cone-graph,gmt:hausdorff-density-upper,gmt:hausdorff-density-outside,gmt:vitali}
Stratify cone directions, apertures, lower-density bounds, and error scales
into countably many sets satisfying the pairwise cone condition.
\end{proof}

\begin{definition}[Weak linear approximation]
\label{def:weak-approximation}
\lean{Preiss.WeaklyLinearlyApproximableAt}
\uses{def:flat}
For each $\eta>0$, choose $a,r_0>0$ such that at every $0<r<r_0$ there
is an $m$-plane $V$, possibly depending on $r$, for which the mass of
$S\cap B(x,r)$ outside the $\eta r$ tube about $x+V$ is less than
$\eta r^m$, and every ball $B(z,\eta r)$ centred at
$z\in(x+V)\cap B(x,r)$ contains at least $a r^m$ of $\Hm|S$.
Both conditions are required (De Lellis 5.4).
\end{definition}

\begin{lemma}[Rectifiable/purely unrectifiable decomposition]
\label{lem:decomposition}
\uses{def:rectifiable}
Every Borel set of finite $\Hm$ measure decomposes, modulo a null set,
into a rectifiable part and a part having zero intersection measure with
every Lipschitz $m$-graph (De Lellis 5.7).
\end{lemma}
\begin{proof}[Planned proof]
\uses{lem:carrier-calculus}
Maximize the measure captured by countable graph unions; take a maximizing
sequence. The source proof works directly with graphs; the image formulation
can use the graph equivalence when needed.
\end{proof}

\begin{theorem}[Marstrand--Mattila set criterion]
\label{thm:mm-set}
\lean{PreissBlueprint.rectifiable_of_weak_approximation}
\uses{def:weak-approximation,def:rectifiable}
For $0<m\leq n$, a Borel set of finite $\Hm$ measure which is weakly
linearly approximable at $\Hm$-almost every point is rectifiable
(De Lellis 5.5; the null set case is immediate).
\end{theorem}
\begin{proof}[Planned proof]
\uses{lem:decomposition,lem:cone-criterion,gmt:hausdorff-density-lower,gmt:hausdorff-density-upper,gmt:hausdorff-density-outside}
Formalize the purely unrectifiable projection argument in De Lellis
Lemmas 5.8--5.9 and Section 5.2. The plane may change with scale; do not
replace this result by the simpler unique-tangent criterion of Theorem 4.8.
\end{proof}

\begin{lemma}[All flat tangents give weak approximation on pieces]
\label{lem:flat-to-approximation}
\uses{def:weak-approximation,def:flat}
On the finite Hausdorff pieces above, positive finite density bounds and
flatness of every tangent imply weak linear approximation almost everywhere.
\end{lemma}
\begin{proof}[Planned proof]
\uses{lem:finite-pieces,thm:matching,thm:doubling,thm:compactness,lem:plane-scaling,gmt:restriction-tangents,gmt:hausdorff-regular}
Transfer tangents to the pieces. Failure of either approximation condition
produces a bad-scale subsequence whose flat limit contradicts portmanteau.
Preserve quantifier order: every tangent is flat, but the plane need not
be independent of the subsequence.
\end{proof}

\begin{theorem}[Marstrand--Mattila measure criterion]
\label{thm:mm-measure}
\lean{PreissBlueprint.rectifiable_of_all_tangents_flat}
\uses{def:positive-density,def:flat,def:rectifiable}
Let $0<m\leq n$ and let $\mu$ be locally finite. If the lower density is
positive and upper density finite almost everywhere, and every tangent at
almost every point is flat, then $\mu$ is $m$-rectifiable (De Lellis 5.1).
\end{theorem}
\begin{proof}[Planned proof]
\uses{lem:absolute-continuity,lem:finite-pieces,lem:flat-to-approximation,thm:mm-set,lem:carrier-calculus}
Apply the set criterion on each piece and assemble their carriers.
\end{proof}

\chapter{Part A: existence of a flat tangent}

\begin{lemma}[Density-to-tangent adapters]
\label{lem:uniform-adapter}
\lean{PreissBlueprint.ae_tangent_uniform,PreissBlueprint.tangent_exists}
\uses{def:positive-density,def:tangent-api,def:uniform-api}
For $m>0$, positive finite existing density almost everywhere implies that
almost everywhere all tangents are uniform and centred at the origin.
Positive finite lower and upper densities give a nonempty tangent set.
\end{lemma}
\begin{proof}[Planned proof]
\uses{lem:density-bridge,thm:uniform-upstream,thm:doubling,gmt:uniform-density-adapter,gmt:doubling-bounds}
Convert the upstream exceptional-set formulation into one a.e. assertion.
Derive support membership and doubling from the two ball growth bounds.
\end{proof}

\begin{lemma}[Closedness of the tangent cone away from zero]
\label{lem:tangent-closed}
\lean{PreissBlueprint.tangent_closed}
\uses{def:tangent-api}
A nonzero regular vague limit of tangent measures at a fixed point is again
a tangent measure of a locally finite measure at that point.
\end{lemma}
\begin{proof}[Planned proof]
\uses{thm:compactness}
Construct a countable determining family of compactly supported tests and
diagonalize the blow-up sequences. Local mass bounds are needed to extend
convergence from the countable family to all tests. A metric on the space
of locally bounded regular measures is a reusable output of this step.
\end{proof}

\begin{lemma}[Recentering tangents almost everywhere]
\label{lem:recenter}
\lean{PreissBlueprint.ae_recenter_tangent}
\uses{def:positive-density,def:flat}
Under the density hypothesis, outside one null set, for every tangent $\nu$,
every $y\in\spt\nu$, and every $r>0$, the measure $D^m_{y,r}\nu$ is
also a tangent at the original point (De Lellis 6.6).
\end{lemma}
\begin{proof}[Planned proof]
\uses{lem:tangent-ray,lem:tangent-closed,thm:matching,lem:plane-scaling,gmt:density-points,gmt:restriction-tangents,gmt:goodsets}
Use the separable test metric and density points of good sets as in Section
6.2. This is not a pointwise identity: the singular vertex of a uniform
nonflat cone shows why the exceptional set matters. The a.e. binder must
precede the universal quantifiers over measures, centres and radii.
\end{proof}

\begin{lemma}[Uniform dimension reduction with blow-up witnesses]
\label{lem:uniform-reduction}
\uses{def:flat,def:uniform-api}
The half-space support reduction for an $m$-uniform measure with ambient
dimension larger than $m$ can be realized by iterated moving-centre limits
of its rescalings. Preserve uniformity and carry the original blow-up
witnesses through the isometric identification of the hyperplane.
\end{lemma}
\begin{proof}[Planned proof]
\uses{thm:compactness,lem:plane-scaling,gmt:uniform-halfspace,gmt:uniform-hyperplane,gmt:hyperplane-transport,gmt:uniform-lower-dim}
Follow De Lellis 3.7--3.9. Audit the candidates
\texttt{SUniformAux.exists\_hyperplane} and
\texttt{SUniformAux.exists\_lower\_dim}. Reusing existence of a smaller
uniform measure without its relationship to the original measure is
insufficient. The analytic-variety alternative (3.11) is not on this path.
\end{proof}

\begin{theorem}[Flat pseudotangent of a uniform measure]
\label{thm:flat-pseudotangent}
\lean{PreissBlueprint.uniform_has_flat_pseudotangent}
\uses{def:flat,def:uniform-api}
For $0<m\leq n$, every $m$-uniform $\nu$ admits a flat limit of
$D^m_{y_i,r_i}\nu$ with $y_i\in\spt\nu$ and $0<r_i\leq1$.
The centres move; no claim is made that this is a tangent at zero
(De Lellis 3.9 and 6.5).
\end{theorem}
\begin{proof}[Planned proof]
\uses{lem:uniform-reduction,lem:plane-scaling,gmt:uniform-full-dimension}
Induct on ambient codimension. At full dimension use
\texttt{mattila\_14\_11}, then flatten the iterated limits by a diagonal argument.
\end{proof}

\begin{theorem}[Part A]
\label{thm:part-a}
\lean{PreissBlueprint.ae_exists_flat_tangent}
\uses{def:positive-density,def:flat}
For $0<m\leq n$, the density hypothesis implies that at almost every point
there exists at least one flat tangent (De Lellis 6.8).
\end{theorem}
\begin{proof}[Planned proof]
\uses{lem:uniform-adapter,lem:recenter,lem:tangent-closed,thm:flat-pseudotangent}
Choose a tangent, apply the pseudotangent theorem, recenter every term,
and use closedness for the nonzero flat limit.
\end{proof}

\chapter{Uniform measures at infinity}

All quantitative results in this chapter use $0<m\leq n$ and unit
uniform measures. De Lellis uses ball mass $\omega_m$ instead of one;
transport formulas by that fixed scalar. In particular a universal energy
gap is false if arbitrary mass scaling is allowed without normalization.

\begin{lemma}[Blow-down compactness and stability]
\label{lem:infinity-compactness}
\uses{def:flat}
For a unit uniform measure, every sequence $r_i\to\infty$ has a
subsequence whose dimensional blow-downs converge to a unit uniform
measure. More generally unit uniform measures form a sequentially compact
class for vague convergence. Limits are nonzero and contain zero in support.
\end{lemma}
\begin{proof}[Planned proof]
\uses{thm:compactness,lem:plane-scaling,gmt:portmanteau,gmt:blowup-balls}
Use uniform ball growth for compact bounds, approximate support points,
and squeeze radii to recover exact ball masses without assuming a priori
that all spheres of an arbitrary limit are null.
\end{proof}

\begin{definition}[Gaussian moments]
\label{def:moments}
\lean{Preiss.gaussianMoment}
For $s>0$, let $I(s)=\int e^{-s|z|^2}\,d\nu(z)$ and
\[
b_{k,s}(u_1,\ldots,u_k)=\frac{(2s)^k}{k!I(s)}
\int\prod_{i=1}^k\langle z,u_i\rangle e^{-s|z|^2}\,d\nu(z).
\]
The present Lean definition is the scalar evaluation. Its packaging as a
symmetric continuous multilinear map is a separate reusable task.
\end{definition}

\begin{lemma}[Radial integrals and Gaussian tails]
\label{lem:radial}
\uses{def:uniform-api,def:moments}
If $\nu$ is unit uniform and $y\in\spt\nu$, radial integration about
$y$ agrees with $\omega_m^{-1}$ times Lebesgue integration on $\R^m$.
Consequently all polynomial Gaussian moments are integrable, $I(s)>0$,
and polynomial Gaussian tails are uniformly controlled under blow-downs.
\end{lemma}
\begin{proof}[Planned proof]
\uses{gmt:uniform-radial}
Start from the existing equality of bounded continuous radial integrals. Extend by approximation and integrable tails to the global Gaussian and Borel radial formulas required here.
\end{proof}
\begin{proof}[Planned proof]
First prove the nonnegative Borel radial formula from the ball distribution;
then derive signed integrable versions. Reuse mathlib Gaussian integration.
Vague convergence alone does not justify passing unbounded tests to limits.
\end{proof}

\begin{lemma}[First Taylor expansion]
\label{lem:taylor-one}
\uses{def:moments}
Uniform measures satisfy the moment bound (7.4) and the finite expansion
(7.5) of De Lellis Lemma 7.6, with constants depending on dimension and
expansion order. In particular, for support points $x$ the remainder of
$\sum_{k=1}^{2q}b_{k,s}(x^k)-\sum_{j=1}^q s^j|x|^{2j}/j!$
is $O((s|x|^2)^{q+1/2})$.
\end{lemma}
\begin{proof}[Planned proof]
\uses{lem:radial}
Apply Taylor estimates to the Gaussian translation identity, with explicit
integrable domination; retain a quantitative remainder before taking limits.
\end{proof}

\begin{lemma}[Moment coefficient expansion]
\label{lem:taylor-two}
\uses{def:moments}
For $k\geq1$, there are symmetric $k$-linear forms $b_k^{(j)}$ such that
for every $q\geq1$,
$b_{k,s}=\sum_{j=1}^q s^j b_k^{(j)}/j!+o(s^q)$ as $s\downarrow0$,
$b_k^{(j)}=0$ for $k>2j$, and
$\sum_{k=1}^{2q}b_k^{(q)}(x^k)=|x|^{2q}$ on the support
(De Lellis 7.7). The zeroth moment is identically one and is handled separately.
\end{lemma}
\begin{proof}[Planned proof]
\uses{lem:taylor-one}
Implement Section 7.3's finite-dimensional polynomial and multilinear
algebra, including the projection onto forms determined on the support.
Do not infer coefficient convergence from boundedness alone. This is a
major proof milestone whose typed tensor API still needs refinement.
\end{proof}

\begin{lemma}[Gaussian moment determination]
\label{lem:moment-determination}
\uses{def:moments}
Two regular measures with the uniform polynomial growth bound are equal
if their Gaussian-weighted polynomial moments agree.
\end{lemma}
\begin{proof}[Planned proof]
\uses{lem:radial}
Use the Gaussian tail control and the approximation argument of Section
7.1. Stone--Weierstrass on a compact set alone does not control polynomial
tails on all of $\R^n$; include the weighted approximation step.
\end{proof}

\begin{theorem}[Unique tangent at infinity]
\label{thm:unique-infinity}
\lean{PreissBlueprint.unique_tangent_at_infinity}
\uses{def:flat}
Every unit uniform measure has exactly one tangent at infinity, which is
unit uniform (De Lellis 6.16). Consequently the entire blow-down family
converges against every compactly supported continuous test as $r\to\infty$.
\end{theorem}
\begin{proof}[Planned proof]
\uses{lem:infinity-compactness,lem:taylor-two,lem:moment-determination,lem:radial}
The expansions give limits of rescaled moments (7.5 in the proposition
numbering). Tail control transfers them to every subsequential limit;
moment determination proves uniqueness, and compactness yields full convergence.
\end{proof}

\begin{corollary}[Conicality at infinity]
\label{cor:conical}
\lean{PreissBlueprint.tangent_at_infinity_conical}
\uses{def:flat}
The unique tangent at infinity of a unit uniform measure is conical
(De Lellis 8.4).
\end{corollary}
\begin{proof}[Planned proof]
\uses{thm:unique-infinity,lem:plane-scaling}
For fixed $a>0$, both a blow-down limit and its $a$-dilation are tangents at
infinity; uniqueness identifies them.
\end{proof}

\begin{lemma}[Conical moment identities]
\label{lem:conical-moments}
\uses{def:flat,def:moments}
For conical unit uniform measures, odd Gaussian moments vanish, each even
moment has just its homogeneous term, the support satisfies the moment
polynomial identities, and the quadratic moment form is positive
semidefinite with trace $m$ (De Lellis 8.6--8.7).
\end{lemma}
\begin{proof}[Planned proof]
\uses{lem:taylor-two,lem:radial}
Combine scaling with the integer-power Taylor expansion, then diagonalize
the quadratic form using existing finite-dimensional spectral theory.
\end{proof}

\begin{theorem}[Gap for conical uniform measures]
\label{thm:gap}
\lean{PreissBlueprint.conical_flatness_gap}
\uses{def:flat}
There is $\varepsilon(m,n)>0$ such that a conical unit uniform $\xi$ is
flat whenever for some $m$-plane $V$,
$\int_{B(0,1)}\dist(z,V)^2\,d\xi(z)\leq\varepsilon(m,n)$
(De Lellis 8.5). This includes the easier dimensions one and two.
\end{theorem}
\begin{proof}[Planned proof]
\uses{lem:conical-moments,lem:infinity-compactness,lem:plane-scaling}
Follow Section 8.3: spectral estimates and support identities rule out
a sequence of nonflat conical uniform measures whose energies tend to zero.
\end{proof}

\begin{lemma}[Moments under flatness at infinity]
\label{lem:flat-infinity-moments}
\uses{def:flat,def:moments}
If the tangent at infinity is the unit flat measure on $V$, then
$b_{2k}^{(k)}(x^{2k})=|P_Vx|^{2k}$ and the restriction of
$b_{2k-1}^{(k)}$ to $V$ vanishes. There is $b\in V^\perp$ with
$|Q_Vx|^2=2\langle b,x\rangle$ on the support, and the support is at
bounded Hausdorff distance from $V$ (De Lellis 9.2--9.4).
\end{lemma}
\begin{proof}[Planned proof]
\uses{thm:unique-infinity,lem:taylor-two,lem:radial,lem:plane-scaling}
Identify leading coefficients by Gaussian moments of blow-down limits.
Use the support polynomial identity at orders one and two.
\end{proof}

\begin{lemma}[Trace rigidity]
\label{lem:trace-rigidity}
\uses{def:moments}
Under flatness at infinity, the coefficients satisfy
$\operatorname{tr}b_2^{(2)}=0$ and
$\operatorname{tr}b_2^{(2)}\geq4\|b_1^{(1)}\|^2/(m+2)$
(De Lellis 9.5). Hence $b_1^{(1)}=0$.
\end{lemma}
\begin{proof}[Planned proof]
\uses{lem:flat-infinity-moments,lem:radial}
Separate the zero-trace calculation, the integral trace formula (9.2 of
the section numbering), the intermediate inequality of Section 9.3, and
the final fourth-order estimate of Section 9.4. These require new proof work;
they are not consequences of elementary compactness.
\end{proof}

\begin{theorem}[Flatness at infinity implies flatness]
\label{thm:infinity-rigidity}
\lean{PreissBlueprint.flat_of_flat_at_infinity}
\uses{def:flat}
If a unit uniform measure has a flat tangent at infinity then the measure
itself is flat (De Lellis 6.19 and 9.1).
\end{theorem}
\begin{proof}[Planned proof]
\uses{lem:trace-rigidity,lem:flat-infinity-moments,lem:plane-scaling}
The vanishing linear coefficient forces the support into $V$; uniformity
then identifies the measure with the unit flat measure on $V$.
\end{proof}

\chapter{Part B and assembly}

\begin{lemma}[Compact space of planes]
\label{lem:grassmannian}
\uses{def:flat}
For $m\leq n$, $m$-planes form a nonempty compact metric space, and
$V\mapsto\dist(z,V)$ is continuous, uniformly in $z$ on compact sets.
\end{lemma}
\begin{proof}[Planned proof]
Represent a plane by its self-adjoint idempotent orthogonal projection of
trace $m$. Prove this is a closed bounded subset of a finite-dimensional
operator space and transport its topology to the plane subtype. Do not
use an unrelated lattice topology on submodules.
\end{proof}

\begin{definition}[Continuous flatness functional]
\label{def:flatness}
\lean{Preiss.flatnessEnergy,Preiss.flatness}
\uses{def:flat}
For a nonnegative compactly supported continuous $\phi$, set
$F_\phi(\nu)=\inf_V\int\phi(z)\dist(z,V)^2\,d\nu(z)$.
The Lean definition is total, but its minimization API requires $m\leq n$.
An infimum over an empty plane type is not geometrically meaningful.
\end{definition}

\begin{lemma}[Continuity of flatness]
\label{lem:flatness-continuous}
\lean{PreissBlueprint.flatness_continuous}
\uses{def:flatness}
For $m\leq n$, the infimum is attained and $F_\phi$ is continuous under
vague convergence of regular measures (De Lellis 6.20).
\end{lemma}
\begin{proof}[Planned proof]
\uses{lem:grassmannian,thm:compactness}
Use compactness of planes, uniform convergence of the integrands on a
fixed compact support, and uniformly bounded local masses.
\end{proof}

\begin{lemma}[First-crossing scales]
\label{lem:first-crossing}
\uses{def:flatness,def:positive-density}
For the normalized blow-up curve at a positive finite density point, a flat
limit and a uniform limit nonflat at infinity yield first-crossing scales
for a fixed positive level of $F_\phi$. After taking a subsequence, the
crossing measures converge and the adjacent scale ratios diverge. The
blow-down of that limit stays below the crossing level.
\end{lemma}
\begin{proof}[Planned proof]
\uses{lem:flatness-continuous,lem:tangent-ray,lem:tangent-closed,thm:compactness,thm:unique-infinity,lem:plane-scaling}
Formalize De Lellis Section 6.4 with $0\leq\phi\leq1$, $\phi=1$ on
$B(0,1)$ and support in $B(0,2)$. Prove continuity in the positive radius,
choose first crossings, and justify the two-scale diagonal limit.
Plain connectedness does not by itself prove the needed large-scale bound.
This node is a technical proof specification, not yet a finalized Lean lemma.
\end{proof}

\begin{theorem}[Part B: one flat tangent forces all tangents flat]
\label{thm:part-b}
\lean{PreissBlueprint.all_tangents_flat_of_one}
\uses{def:flat,def:positive-density}
Let $0<m\leq n$ and let $x$ be a positive finite density point of locally
finite $\mu$. If every tangent at $x$ is $m$-uniform and one is flat, then
every tangent at $x$ is flat (De Lellis 6.10).
\end{theorem}
\begin{proof}[Planned proof]
\uses{lem:first-crossing,thm:gap,cor:conical,thm:infinity-rigidity}
Normalize the common dimensional tangent mass. A nonflat tangent would
give a first-crossing limit which is nonflat but has a sufficiently flat
conical tangent at infinity. The gap and rigidity theorems contradict this.
Do not assert that every uniform measure is flat: the Kowalski--Preiss
three-dimensional cone in $\R^4$ is a counterexample.
\end{proof}

\begin{lemma}[Dimension zero]
\label{lem:zero-dimension}
\lean{PreissBlueprint.preiss_zero}
\uses{def:positive-density,def:rectifiable}
In dimension zero the density hypothesis implies rectifiability.
\end{lemma}
\begin{proof}[Planned proof]
Continuity from above of the finite measure on a bounded ball identifies
the small-ball limit with $\mu\{x\}$. The measure is carried by its atoms;
local finiteness and a countable compact exhaustion make these countable.
Zero-dimensional Hausdorff measure has no nonempty null sets.
\end{proof}

\begin{lemma}[Dimensions above the ambient dimension]
\label{lem:excess-dimension}
\lean{PreissBlueprint.measure_zero_of_dimension_gt,PreissBlueprint.zero_isRectifiableMeasure}
\uses{def:positive-density,def:rectifiable}
If $m>n$, the density hypothesis forces $\mu=0$. The zero measure is
rectifiable in every dimension.
\end{lemma}
\begin{proof}[Planned proof]
\uses{lem:absolute-continuity}
Hausdorff measure in dimension greater than $n$ vanishes on $\R^n$.
For the zero measure use the empty countable carrier.
\end{proof}

\begin{theorem}[Preiss]
\label{thm:preiss}
\lean{PreissBlueprint.preiss}
\uses{def:positive-density,def:rectifiable}
Let $m,n\in\N$ and let $\mu$ be a locally finite Borel measure on $\R^n$.
If $0<\Theta_*^m(\mu,x)=\Theta^{*m}(\mu,x)<\infty$ for
$\mu$-almost every $x$, then $\mu$ is an $m$-rectifiable measure.
\end{theorem}
\begin{proof}[Assembly, still dependent on unproved obligations]
\uses{lem:uniform-adapter,thm:part-a,thm:part-b,thm:mm-measure,lem:zero-dimension,lem:excess-dimension}
For $0<m\leq n$, intersect the finitely many full-measure sets, apply
Parts A and B pointwise, and use the measure criterion. Handle $m=0$ and
$m>n$ separately. The Lean assembly already expresses this dependency
structure, but inherits the unproved obligations' axioms.
\end{proof}

\chapter{Implementation order and review gates}
First stabilize the definitions, prove the density bridges, carrier calculus,
plane/scaling identities, and the two upstream adapters. This produces useful
standalone contributions before any deep rigidity proof is attempted.

Next develop two independent branches: the rectifiability criteria of
Chapters 4--5, and the tangent/pseudotangent calculus for Part A. In the
second branch build the separable vague test metric and retain moving-centre
witnesses in dimension reduction.

The long branch is the moment calculus: radial integration, Gaussian tails,
the two Taylor expansions, moment determination, and uniqueness at infinity.
From these build conical identities and the flatness gap, and separately
the trace identities leading to rigidity. Only then discharge the
first-crossing argument and Part B.

Keep the full area formula, the converse rectifiability/density theorem,
the integer-dimension Marstrand theorem, analytic support varieties, and
the stronger quantitative density-ratio version of Preiss as optional
extensions. They are not silently prerequisites of the present assembly.

Before promoting any theorem out of \texttt{PreissBlueprint}, remove its
\texttt{sorry}, check its actual axiom dependencies, and update the blueprint
status. Before changing upstream revisions, repeat the source audit and
build both targets. The graph is a living proof plan: inferred edges and
technical sublemma boundaries must be revised when actual proofs require it.
